\hypertarget{ituxe9ration}{%
\section{Itération}\label{ituxe9ration}}

\hypertarget{ruxe9affectation}{%
\subsection{Réaffectation}\label{ruxe9affectation}}

Comme vous l\textquotesingle avez peut-être déjà découvert, il est légal
d\textquotesingle effectuer plus d\textquotesingle une affectation sur
la même variable. Une nouvelle affectation fait en sorte
qu\textquotesingle une variable existante se réfère à une nouvelle
valeur (et cesse de se référer à l\textquotesingle ancienne).

\begin{Shaded}
\begin{Highlighting}[]
\OperatorTok{\textgreater{}\textgreater{}\textgreater{}}\NormalTok{ x }\OperatorTok{=} \DecValTok{5}
\OperatorTok{\textgreater{}\textgreater{}\textgreater{}}\NormalTok{ x}
\DecValTok{5}
\OperatorTok{\textgreater{}\textgreater{}\textgreater{}}\NormalTok{ x }\OperatorTok{=} \DecValTok{7}
\OperatorTok{\textgreater{}\textgreater{}\textgreater{}}\NormalTok{ x}
\DecValTok{7}
\end{Highlighting}
\end{Shaded}

La première fois que nous affichons \texttt{x}, sa valeur est 5 ; la
deuxième fois, sa valeur est 7.

L\textquotesingle une des sources de confusion les plus courantes lors
de l\textquotesingle apprentissage de la programmation est la différence
entre l\textquotesingle affectation et l\textquotesingle égalité. En
mathématiques, le signe égal indique que deux éléments sont égaux et le
resteront. En Python, une instruction d\textquotesingle affectation
utilise le signe égal (\texttt{=}), mais ce n\textquotesingle est pas
une affirmation d\textquotesingle égalité.

Par exemple, si vous écrivez \texttt{a\ =\ b}, vous affirmez que
\texttt{a} et \texttt{b} ont actuellement la même valeur, mais ils ne
sont pas obligés de rester égaux indéfiniment :

\begin{Shaded}
\begin{Highlighting}[]
\NormalTok{a }\OperatorTok{=} \DecValTok{5}
\NormalTok{b }\OperatorTok{=}\NormalTok{ a    }\CommentTok{\# a et b sont maintenant égaux}
\NormalTok{a }\OperatorTok{=} \DecValTok{3}    \CommentTok{\# a et b ne sont plus égaux}
\end{Highlighting}
\end{Shaded}

\hypertarget{mise-uxe0-jour-de-variables}{%
\subsection{Mise à jour de
variables}\label{mise-uxe0-jour-de-variables}}

Un des types de réaffectation les plus courants est la mise à jour
(\emph{update}), où la nouvelle valeur de la variable dépend de
l\textquotesingle ancienne.

\begin{Shaded}
\begin{Highlighting}[]
\NormalTok{x }\OperatorTok{=}\NormalTok{ x }\OperatorTok{+} \DecValTok{1}
\end{Highlighting}
\end{Shaded}

Cela signifie « récupère la valeur actuelle de \texttt{x}, ajoute 1, et
met ensuite à jour \texttt{x} avec la nouvelle valeur. »

Si vous essayez de mettre à jour une variable qui
n\textquotesingle existe pas, vous obtenez une erreur, car Python évalue
le côté droit avant d\textquotesingle affecter la valeur à \texttt{x} :

\begin{Shaded}
\begin{Highlighting}[]
\OperatorTok{\textgreater{}\textgreater{}\textgreater{}}\NormalTok{ x }\OperatorTok{=}\NormalTok{ x }\OperatorTok{+} \DecValTok{1}
\PreprocessorTok{NameError}\NormalTok{: name }\StringTok{\textquotesingle{}x\textquotesingle{}} \KeywordTok{is} \KeywordTok{not}\NormalTok{ defined}
\end{Highlighting}
\end{Shaded}

Avant de pouvoir mettre à jour une variable, vous devez
l\textquotesingle initialiser (\emph{initialize}), généralement par une
affectation simple :

\begin{Shaded}
\begin{Highlighting}[]
\OperatorTok{\textgreater{}\textgreater{}\textgreater{}}\NormalTok{ x }\OperatorTok{=} \DecValTok{0}
\OperatorTok{\textgreater{}\textgreater{}\textgreater{}}\NormalTok{ x }\OperatorTok{=}\NormalTok{ x }\OperatorTok{+} \DecValTok{1}
\end{Highlighting}
\end{Shaded}

La mise à jour d\textquotesingle une variable par
l\textquotesingle ajout de 1 s\textquotesingle appelle une
incrémentation (\emph{increment}) ; la soustraction de 1
s\textquotesingle appelle une décrémentation (\emph{decrement}).

\hypertarget{linstruction-while}{%
\subsection{\texorpdfstring{L\textquotesingle instruction
\texttt{while}}{L\textquotesingle instruction while}}\label{linstruction-while}}

Les ordinateurs sont souvent utilisés pour automatiser des tâches
répétitives. Répéter des tâches identiques ou similaires sans faire
d\textquotesingle erreurs est une chose que les ordinateurs font bien et
que les humains font mal. En programmation, la répétition est aussi
appelée itération (\emph{iteration}).

Python fournit une instruction \texttt{while} pour faciliter
l\textquotesingle itération. Voici une version de
\texttt{compte\_a\_rebours} qui utilise une instruction \texttt{while} :

\begin{Shaded}
\begin{Highlighting}[]
\KeywordTok{def}\NormalTok{ compte\_a\_rebours(n):}
    \ControlFlowTok{while}\NormalTok{ n }\OperatorTok{\textgreater{}} \DecValTok{0}\NormalTok{:}
        \BuiltInTok{print}\NormalTok{(n)}
\NormalTok{        n }\OperatorTok{=}\NormalTok{ n }\OperatorTok{{-}} \DecValTok{1}
    \BuiltInTok{print}\NormalTok{(}\StringTok{\textquotesingle{}Décollage !\textquotesingle{}}\NormalTok{)}
\end{Highlighting}
\end{Shaded}

Vous pouvez lire l\textquotesingle instruction \texttt{while} presque
comme de l\textquotesingle anglais. Elle signifie : « Tant que
\texttt{n} est strictement supérieur à 0, affiche la valeur de
\texttt{n} et diminue la valeur de \texttt{n} de 1. Une fois que tu as
terminé, affiche le mot Décollage ! »

Plus formellement, voici le flux d\textquotesingle exécution pour une
instruction \texttt{while} :

\begin{enumerate}
\def\labelenumi{\arabic{enumi}.}
\tightlist
\item
  Déterminez si la condition est vraie ou fausse.
\item
  Si la condition est fausse, quittez l\textquotesingle instruction
  \texttt{while} et continuez l\textquotesingle exécution à la prochaine
  instruction.
\item
  Si la condition est vraie, exécutez le corps de la boucle et retournez
  à l\textquotesingle étape 1.
\end{enumerate}

Ce type de flux est appelé une boucle (\emph{loop}). Si la condition est
fausse dès le départ, le corps de la boucle ne s\textquotesingle exécute
jamais.

Le corps de la boucle doit changer la valeur d\textquotesingle une ou
plusieurs variables pour qu\textquotesingle à un moment donné, la
condition devienne fausse et que la boucle se termine. Si la condition
ne devient jamais fausse, la boucle se répétera pour toujours, ce
qu\textquotesingle on appelle une boucle infinie (\emph{infinite loop}).

\hypertarget{linstruction-break}{%
\subsection{\texorpdfstring{L\textquotesingle instruction
\texttt{break}}{L\textquotesingle instruction break}}\label{linstruction-break}}

Parfois, vous ne savez pas qu\textquotesingle il est temps de terminer
une boucle avant d\textquotesingle être au milieu du corps de la boucle.
Dans ce cas, vous pouvez utiliser l\textquotesingle instruction
\texttt{break} pour sortir de la boucle immédiatement.

Par exemple, supposons que vous souhaitiez demander une saisie à
l\textquotesingle utilisateur jusqu\textquotesingle à ce
qu\textquotesingle il tape
\texttt{\textquotesingle{}fin\textquotesingle{}}. Vous pourriez écrire :

\begin{Shaded}
\begin{Highlighting}[]
\ControlFlowTok{while} \VariableTok{True}\NormalTok{:}
\NormalTok{    ligne }\OperatorTok{=} \BuiltInTok{input}\NormalTok{(}\StringTok{\textquotesingle{}\textgreater{} \textquotesingle{}}\NormalTok{)}
    \ControlFlowTok{if}\NormalTok{ ligne }\OperatorTok{==} \StringTok{\textquotesingle{}fin\textquotesingle{}}\NormalTok{:}
        \ControlFlowTok{break}
    \BuiltInTok{print}\NormalTok{(ligne)}
\BuiltInTok{print}\NormalTok{(}\StringTok{\textquotesingle{}Terminé !\textquotesingle{}}\NormalTok{)}
\end{Highlighting}
\end{Shaded}

La condition de la boucle est \texttt{True}, ce qui est toujours vrai,
donc la boucle s\textquotesingle exécutera jusqu\textquotesingle à ce
qu\textquotesingle elle rencontre l\textquotesingle instruction
\texttt{break}. À chaque itération, elle invite
l\textquotesingle utilisateur avec un chevron. Si
l\textquotesingle utilisateur tape \texttt{fin},
l\textquotesingle instruction \texttt{break} quitte la boucle. Sinon, le
programme affiche ce que l\textquotesingle utilisateur a tapé et
retourne au début de la boucle.

\hypertarget{racines-carruxe9es}{%
\subsection{Racines carrées}\label{racines-carruxe9es}}

Les boucles sont souvent utilisées dans les programmes qui calculent des
résultats numériques en commençant par une réponse approximative et en
l\textquotesingle améliorant de manière itérative.

Par exemple, l\textquotesingle une des façons de calculer les racines
carrées est la méthode de Newton. Supposons que vous souhaitiez
connaître la racine carrée de \$a\$. Si vous commencez avec une
estimation quelconque, \$x\$, vous pouvez calculer une meilleure
estimation avec la formule suivante :

\begin{quote}
\$y = frac\{x + a / x\}\{2\}\$
\end{quote}

Par exemple, si \$a\$ vaut 4 et \$x\$ vaut 3 :

\begin{Shaded}
\begin{Highlighting}[]
\OperatorTok{\textgreater{}\textgreater{}\textgreater{}}\NormalTok{ a }\OperatorTok{=} \DecValTok{4}
\OperatorTok{\textgreater{}\textgreater{}\textgreater{}}\NormalTok{ x }\OperatorTok{=} \DecValTok{3}
\OperatorTok{\textgreater{}\textgreater{}\textgreater{}}\NormalTok{ y }\OperatorTok{=}\NormalTok{ (x }\OperatorTok{+}\NormalTok{ a}\OperatorTok{/}\NormalTok{x) }\OperatorTok{/} \DecValTok{2}
\OperatorTok{\textgreater{}\textgreater{}\textgreater{}}\NormalTok{ y}
\FloatTok{2.1666666666666665}
\end{Highlighting}
\end{Shaded}

Le résultat est plus proche de la bonne réponse (\$sqrt\{4\} = 2\$). Si
nous répétons le processus avec la nouvelle estimation, le résultat
s\textquotesingle améliore encore :

\begin{Shaded}
\begin{Highlighting}[]
\OperatorTok{\textgreater{}\textgreater{}\textgreater{}}\NormalTok{ x }\OperatorTok{=}\NormalTok{ y}
\OperatorTok{\textgreater{}\textgreater{}\textgreater{}}\NormalTok{ y }\OperatorTok{=}\NormalTok{ (x }\OperatorTok{+}\NormalTok{ a}\OperatorTok{/}\NormalTok{x) }\OperatorTok{/} \DecValTok{2}
\OperatorTok{\textgreater{}\textgreater{}\textgreater{}}\NormalTok{ y}
\FloatTok{2.0064102564102564}
\end{Highlighting}
\end{Shaded}

Lorsque \$y == x\$, vous pouvez arrêter le calcul. Vous pouvez exprimer
cela dans une boucle :

\begin{Shaded}
\begin{Highlighting}[]
\ControlFlowTok{while} \VariableTok{True}\NormalTok{:}
    \BuiltInTok{print}\NormalTok{(x)}
\NormalTok{    y }\OperatorTok{=}\NormalTok{ (x }\OperatorTok{+}\NormalTok{ a}\OperatorTok{/}\NormalTok{x) }\OperatorTok{/} \DecValTok{2}
    \ControlFlowTok{if}\NormalTok{ y }\OperatorTok{==}\NormalTok{ x:}
        \ControlFlowTok{break}
\NormalTok{    x }\OperatorTok{=}\NormalTok{ y}
\end{Highlighting}
\end{Shaded}

Pour la plupart des valeurs de \$a\$, cela fonctionne bien, mais de
manière générale, tester l\textquotesingle égalité stricte des nombres à
virgule flottante est dangereux. Plutôt que de vérifier si \texttt{x} et
\texttt{y} sont exactement égaux, il est plus sûr
d\textquotesingle utiliser la fonction intégrée \texttt{abs} pour
calculer la valeur absolue de leur différence :

\begin{Shaded}
\begin{Highlighting}[]
\ControlFlowTok{if} \BuiltInTok{abs}\NormalTok{(y }\OperatorTok{{-}}\NormalTok{ x) }\OperatorTok{\textless{}}\NormalTok{ epsilon:}
    \ControlFlowTok{break}
\end{Highlighting}
\end{Shaded}

où \texttt{epsilon} a une valeur petite comme \texttt{0.0000001} qui
détermine la précision requise.

\hypertarget{algorithmes}{%
\subsection{Algorithmes}\label{algorithmes}}

La méthode de Newton est un exemple d\textquotesingle algorithme
(\emph{algorithm}) : c\textquotesingle est un processus mécanique pour
résoudre une catégorie de problèmes (dans ce cas, calculer des racines
carrées).

Pour comprendre ce qu\textquotesingle est un algorithme, il peut être
utile de commencer par ce qui n\textquotesingle en est pas un.
L\textquotesingle arithmétique de base que vous avez apprise
n\textquotesingle est généralement pas algorithmique. Mais si vous avez
appris la division longue, vous avez appris un algorithme. Un algorithme
est une suite d\textquotesingle étapes systématiques, sans nécessiter
d\textquotesingle intelligence ni d\textquotesingle intuition pour les
exécuter.

\hypertarget{duxe9bogage}{%
\subsection{Débogage}\label{duxe9bogage}}

À mesure que vous commencez à écrire des programmes plus longs, vous
pourriez vous retrouver à passer plus de temps à déboguer. Réduire la
quantité de code que vous ajoutez avant de tester (développement
incrémental) peut aider.

Une autre technique utile est la division par deux de la recherche
(\emph{bisection}). Si vous avez 100 lignes de code et
qu\textquotesingle il y a une erreur, vous pouvez placer une instruction
\texttt{print} au milieu pour vérifier si l\textquotesingle état du
programme est correct. S\textquotesingle il l\textquotesingle est,
l\textquotesingle erreur se trouve dans la seconde moitié ; sinon, elle
est dans la première. En coupant à chaque fois la zone de recherche en
deux, vous trouverez plus vite le problème.

\hypertarget{glossaire}{%
\subsection{Glossaire}\label{glossaire}}

réaffectation reassignment Le fait d\textquotesingle assigner une
nouvelle valeur à une variable qui existait déjà.

mise à jour update Une affectation où la nouvelle valeur de la variable
dépend de l\textquotesingle ancienne.

initialisation initialization Une affectation qui donne une valeur
initiale à une variable qui sera mise à jour par la suite.

incrémenter increment Une mise à jour qui augmente la valeur
d\textquotesingle une variable (souvent de 1).

décrémenter decrement Une mise à jour qui diminue la valeur
d\textquotesingle une variable.

itération iteration L\textquotesingle exécution répétée
d\textquotesingle un ensemble d\textquotesingle instructions en
utilisant soit une fonction récursive, soit une boucle.

boucle infinie infinite loop Une boucle dont la condition de terminaison
n\textquotesingle est jamais satisfaite.

algorithme algorithm Un processus général pour résoudre une catégorie de
problèmes.

\hypertarget{exercices}{%
\subsection{Exercices}\label{exercices}}

\textbf{Exercice 1}

Copiez la boucle de la section sur les racines carrées et encapsulez-la
dans une fonction appelée \texttt{ma\_racine(a)}.

Pour tester votre fonction, écrivez une fonction nommée
\texttt{tester\_racine\_carree} qui affiche un tableau semblable à
celui-ci :

\begin{Shaded}
\begin{Highlighting}[]
\NormalTok{a   ma\_racine(a)      math.sqrt(a)   diff}
\NormalTok{{-}   {-}{-}{-}{-}{-}{-}{-}{-}{-}{-}{-}{-}      {-}{-}{-}{-}{-}{-}{-}{-}{-}{-}{-}{-}   {-}{-}{-}{-}}
\NormalTok{1.0 1.0               1.0            0.0}
\NormalTok{2.0 1.41421356237     1.41421356237  2.22044604925e{-}16}
\NormalTok{3.0 1.73205081001     1.73205081001  0.0}
\NormalTok{4.0 2.0               2.0            0.0}
\end{Highlighting}
\end{Shaded}

La première colonne correspond aux nombres de 1 à 9. La deuxième colonne
est le résultat calculé par votre fonction \texttt{ma\_racine}. La
troisième est le résultat de \texttt{math.sqrt}, et la quatrième colonne
est la différence absolue entre les deux.

\textbf{Exercice 2}

La fonction intégrée \texttt{eval} prend une chaîne de caractères et
l\textquotesingle évalue à l\textquotesingle aide de
l\textquotesingle interpréteur Python. Par exemple :

\begin{Shaded}
\begin{Highlighting}[]
\OperatorTok{\textgreater{}\textgreater{}\textgreater{}} \BuiltInTok{eval}\NormalTok{(}\StringTok{\textquotesingle{}1 + 2 * 3\textquotesingle{}}\NormalTok{)}
\DecValTok{7}
\OperatorTok{\textgreater{}\textgreater{}\textgreater{}} \ImportTok{import}\NormalTok{ math}
\OperatorTok{\textgreater{}\textgreater{}\textgreater{}} \BuiltInTok{eval}\NormalTok{(}\StringTok{\textquotesingle{}math.sqrt(5)\textquotesingle{}}\NormalTok{)}
\FloatTok{2.2360679774997898}
\OperatorTok{\textgreater{}\textgreater{}\textgreater{}} \BuiltInTok{eval}\NormalTok{(}\StringTok{\textquotesingle{}type(math.pi)\textquotesingle{}}\NormalTok{)}
\OperatorTok{\textless{}}\KeywordTok{class} \StringTok{\textquotesingle{}float\textquotesingle{}}\OperatorTok{\textgreater{}}
\end{Highlighting}
\end{Shaded}

Écrivez une fonction appelée \texttt{eval\_loop} qui invite
itérativement l\textquotesingle utilisateur à entrer une saisie, prend
la saisie, et l\textquotesingle évalue en utilisant \texttt{eval}, pour
finalement afficher le résultat.

Elle devrait continuer ainsi jusqu\textquotesingle à ce que
l\textquotesingle utilisateur tape
\texttt{\textquotesingle{}done\textquotesingle{}}, et elle devrait
ensuite retourner la valeur de la dernière expression évaluée.

\textbf{Exercice 3}

Le brillant mathématicien Srinivasa Ramanujan a découvert une série
infinie permettant de calculer une approximation de \$pi\$ :

\$\$ frac\{1\}\{pi\} = frac\{2sqrt\{2\}\}\{9801\}
s\href{}{um}\{k=0\}\^{}\{infty\} frac\{(4k)!(1103 + 26390k)\}\{(k!)\^{}4
396\^{}\{4k\}\} \$\$

Écrivez une fonction appelée \texttt{estimer\_pi} qui utilise cette
formule pour calculer et retourner une estimation de \$pi\$. Elle doit
utiliser une boucle \texttt{while} pour calculer les termes de la somme
jusqu\textquotesingle à ce que le dernier terme calculé soit strictement
inférieur à \texttt{1e-15} (soit la notation scientifique de Python pour
\$10\^{}\{-15\}\$). Vous pourrez vérifier votre résultat en le comparant
à \texttt{math.pi}.
